\section{Correctness}\label{sec:correctness}

Throughout, $(\mathit{pk},\mathit{sk})$ is honestly generated,
$\mathit{ct}=\textsf{Enc}(\mathit{pk},s)$, and the notation of
Algorithms~\ref{alg:keygen}--\ref{alg:decaps} is in force. Write
$\theta=\varphi\tau\hat\varphi\in\End(\EC)$ and $\psi=2\theta-tdf$.

\begin{lemma}[the lollipop endomorphism]\label{lem:psi}
$\Tr\psi=0$ and $\deg\psi=7d^2f^2$. Moreover $\psi S=ds_0S$ and
$\psi T=-ds_0T$, where $s_0=fs_D\bmod 2^E$.
\end{lemma}

\begin{proof}
$\Tr\theta=\varphi(\tau+\hat\tau)\hat\varphi=\varphi[tf]\hat\varphi=tdf$, and
$\deg\theta=d^2\deg\tau=2d^2f^2$. Hence $\Tr\psi=2\Tr\theta-2tdf=0$ and
\[
\deg\psi=(2\theta-tdf)(2\hat\theta-tdf)=4\deg\theta-2tdf\Tr\theta+t^2d^2f^2=8d^2f^2-2d^2f^2+d^2f^2=7d^2f^2 .
\]
For the eigenvalues, $\hat\varphi S=a\hat\varphi\varphi P_A=adP_A$, and
$\tau P_A=f\mu P_A$, so $\theta S=adf\mu\,\varphi P_A=df\mu S$. Likewise
$\theta T=df\nu T$. Then $\psi S=df(2\mu-t)S=dfs_DS$, and
$\psi T=df(2\nu-t)T=-dfs_DT$.
\end{proof}

The masks $a,b$ drop out because $S$ and $T$ are multiples of eigenvectors.
This is where the trapdoor is used: only $f$ (through $s_0$) and the Kani
scalars are needed, never $\varphi$ or the masks.

\begin{proposition}[the decryption isogeny]\label{prop:kani}
Let $\gamma=\psi+c$. Then $\deg\gamma+c'^2=2^E$, and the subgroup
$\langle(c'S,-\lambda_SS),(c'T,-\lambda_TT)\rangle$ of $\EC^2$ is the kernel of
the $(2^E,2^E)$-isogeny $\Phi=\begin{psmallmatrix}c'&-\hat\gamma\\ \gamma&c'\end{psmallmatrix}$.
\end{proposition}

\begin{proof}
$\deg\gamma=\deg\psi+c\Tr\psi+c^2=7d^2f^2+c^2=2^E-c'^2$ by the choice of
$c,c'$. By Lemma~\ref{lem:psi}, $\gamma S=(ds_0+c)S=\lambda_SS$ and
$\gamma T=(c-ds_0)T=\lambda_TT$. The points $S,T$ form a basis of $\EC[2^E]$,
because $\varphi$ has odd degree and $a,b$ are odd. So the two generators
are $(c'R,-\gamma R)$ for $R=S,T$, and the kernel statement is Kani's lemma
(\S\ref{sec:prelim}).
\end{proof}

From $\Phi(X,0)=(c'X,\gamma X)$, decapsulation obtains $\gamma$, and hence
$\psi=\gamma-c$, on $\EC[d]$. Restricted to $\EC[d]$ we have
$\psi=2\varphi\tau\hat\varphi$, since $d\mid tdf$.

\begin{proposition}[the recovered kernel is always cyclic]\label{prop:cyclic}
$\psi(\EC[d])$ is cyclic of order $d$, and the kernel of $\psi$ restricted to
$\EC[d]$ is exactly $\ker\hat\varphi$.
\end{proposition}

\begin{proof}
We have $\hat\varphi(\EC[d])=\ker\varphi=\langle K\rangle$, which is cyclic of
order $d$. The endomorphism $\tau$ has degree $2f^2$, which is prime to $3$,
so $\tau\langle K\rangle$ is again cyclic of order $d$. The map $\varphi$ is
injective on $\tau\langle K\rangle$ if and only if
$\tau\langle K\rangle\cap\langle K\rangle=0$, i.e.\ if and only if the two
subgroups have different subgroups of order~$3$. That fails only if the line
$L=\langle 3^{F-1}K\rangle\subset\EA[3]$ is an eigenline of $\tau$ modulo~3.
The characteristic polynomial of $\tau$ is $x^2-tfx+2f^2$, whose
discriminant is $f^2(t^2-8)=-7f^2\equiv 2f^2\pmod 3$. Since $3\nmid f$, this is
a non-square modulo~3. So $\tau$ has no eigenline in $\EA[3]$, and $\varphi$
maps $\tau\langle K\rangle$ isomorphically onto a cyclic group of order $d$.
Thus $\psi(\EC[d])=2\varphi\tau\langle K\rangle$ has order $d$, its kernel
has order $d^2/d=d$, and that kernel contains $\ker\hat\varphi$, which also
has order~$d$.
\end{proof}

\begin{remark}\label{rem:noncyclic}
The design note~\cite{EigenSetaNote} reported a ``non-cyclic $\psi'$'' corner
case in one of eight trials of its toy prototype, a variant with a different
$\psi'$ and very small $d=3^5$. Proposition~\ref{prop:cyclic} shows that the
case cannot occur in version~2. The argument needs only that $-7$ is a
non-square modulo~3 and $3\nmid\deg\tau$. The cause of the toy observation
was not isolated.
\end{remark}

\begin{proposition}[recovery of $s$]\label{prop:s}
Let $G$ be a generator of $\psi(\EC[d])$ and $H=kG$ the image of the other
basis point. Then $\ker\hat\varphi=\langle X_h-kX_g\rangle$. If $K_g$ is the image
under $\hat\varphi$ of a basis point and has order $d$, then $K_g=u(P_3+sQ_3)$
for a unit $u$, and
$s=\log e(P_3,K_g)/\log e(K_g,Q_3)$ in $\mu_d$.
\end{proposition}

\begin{proof}
$\psi(X_h-kX_g)=H-kG=0$, and $X_h-kX_g$ has order $d$, so by
Proposition~\ref{prop:cyclic} it generates $\ker\hat\varphi$. The image of
$\EC[d]$ under $\hat\varphi$ is $\ker\varphi=\langle P_3+sQ_3\rangle$. With
$\zeta=e(P_3,Q_3)$, which has order $d$, we get $e(K_g,Q_3)=\zeta^u$ and
$e(P_3,K_g)=\zeta^{us}$. The codomain of $\hat\varphi$ is $\EA$ up to
isomorphism, and isomorphisms preserve these subgroups. The implementation
maps the codomain to the model with coefficient $A$, using $x\mapsto-x$ when
the computed coefficient is $-A$.
\end{proof}

\begin{proposition}[fixed schedule]\label{prop:schedule}
In the basis $(S,0),(T,0),(0,S),(0,T)$ of $\EC^2[2^E]$, the coordinates of
$\ker\Phi$ are the columns of
$\begin{psmallmatrix}c'&0\\0&c'\\-\lambda_S&0\\0&-\lambda_T\end{psmallmatrix}$
modulo $2^E$, which depend on $\mathit{sk}$ only. Whether each $(2,2)$-step of
the chain is a product step, a gluing, a Richelot step or a splitting is
decided by linear conditions on these coordinates modulo powers of~2. So the
schedule is the same for every ciphertext under a given key.
\end{proposition}

\begin{proof}[Proof sketch]
The kernel of the $j$-th step is $2^{E-j}$ times the image of $\ker\Phi$ under
the first $j-1$ steps. On a product $E_1\times E_2$, the step is a product
step exactly when its 2-torsion kernel is a direct product
$\langle(P,0),(0,Q)\rangle$, and a graph step when it is the graph of the
identity of $\EC[2]$ (on $\EC\times\EC$) or of an isomorphism (after the
split). Both are statements about the coordinates. On a Jacobian, the step
is a splitting exactly when the codomain is a product, which happens at a
fixed index determined by the same 2-adic data. Key generation imposes
$\max(v_2(c'\mp\lambda_S))=2$ and the gluing test of line~15, which fixes the
schedule as graph, product, glue, Richelot${}^{E-6}$, split, product, graph.
The self-test checks it once.
\end{proof}

\paragraph{Correctness error.} By the propositions above, decapsulation of an
honest ciphertext returns $H_K(s,\mathit{ct})$ unless an intermediate formula
meets a degenerate input. Examples are a Mumford divisor not in general
position, a zero denominator in a Richelot or splitting formula, or a Miller
line through the second argument. No branch handles these cases. The output
is then wrong but computed in the same time, and the FO check rejects it.
Heuristically each such event has probability $O(1/p)$ per step, giving
$\delta=O(E/p)\approx2^{-1330}$ at level~1. Key generation excludes the one
structured degenerate case, a scalar $\equiv-1$ modulo the point order, which
the x-only ladder with $y$-recovery cannot handle. The measured rates are in
Table~\ref{tab:failures}: no failure occurred in \FailTotalTrials\
decapsulations.
